\documentclass[10pt,letterpaper]{article}
\usepackage[T1]{fontenc}
\usepackage[utf8]{inputenc}
\usepackage{lmodern}
\usepackage{microtype}
\usepackage[margin=1in]{geometry}
\usepackage{array}
\usepackage{booktabs}
\usepackage{longtable}
\usepackage{enumitem}
\usepackage[hidelinks]{hyperref}
\setlength{\parindent}{0pt}
\setlength{\parskip}{0.5em}
\renewcommand{\arraystretch}{1.15}
\setlist[itemize]{leftmargin=*,itemsep=0.25em,topsep=0.3em}

\title{\textbf{Testing Equipment Attachment}\\900 MHz Gilbert-Cell Downconversion Mixer}
\author{Sangeon Jeon, Yunho Cho}
\date{November 2026}

\begin{document}
\maketitle

\section{Facility Inventory}

This separate attachment supports Chipalooza Challenge \#4 Rule 6.6.
All 21 instrument and assembly-tool entries on The Hive's public electronics-benchtop inventory are included below [1].
The proposed uses are our measurement plan; the inventory does not establish installed options, calibration status, booking, or simultaneous availability.
In particular, the vector signal generator is listed without a model number, and the 33500B designation identifies a product family rather than a specific configuration.

\small
\begin{longtable}{@{}>{\raggedright\arraybackslash}p{0.43\textwidth}>{\raggedright\arraybackslash}p{0.53\textwidth}@{}}
\toprule
\textbf{Listed equipment} & \textbf{Proposed use} \\
\midrule
\endfirsthead
\toprule
\textbf{Listed equipment (continued)} & \textbf{Proposed use} \\
\midrule
\endhead
\midrule
\multicolumn{2}{r}{Continued on next page} \\
\endfoot
\bottomrule
\endlastfoot

\multicolumn{2}{@{}l}{\textbf{RF, IF, and impedance measurements}} \\
Keysight E5063A & RF input return loss and fixture characterization, conditional on the installed frequency option and calibration kit. \\
Keysight N9000B & IF fundamental and intermodulation spectra; RF/LO leakage measurements with appropriate coupling, attenuation, and calibration. \\
Keysight vector signal generator & Candidate RF or LO stimulus; exact model, range, options, and available source count to be confirmed. \\
Keysight N1913A & RF/LO power calibration with a compatible sensor covering the required frequency and power range. \\
Keysight E4980-AL (inventory spelling) & Passive-component and fixture impedance checks within the instrument's operating range. \\
\addlinespace

\multicolumn{2}{@{}l}{\textbf{Waveform generation and observation}} \\
Tektronix MSO58 & IF waveform, settling, startup, and control-signal checks using suitable probes. \\
Keysight DSOX3014A & IF waveform and transient checks. \\
Keysight DSOX4024A & IF waveform and transient checks. \\
Keysight DSOX6002A & IF waveform and transient checks; direct RF probing depends on installed bandwidth and probe loading. \\
Tektronix AFG31152 & IF/baseband stimulus and buffer or fixture checks; not a direct 900 MHz source [3]. \\
Keysight 33500B family & IF/baseband stimulus; not a direct 900 MHz source [4]. \\
\addlinespace

\multicolumn{2}{@{}l}{\textbf{DC supply, bias, and operating-point measurements}} \\
Keithley 2450 & Controlled DC sourcing and current/voltage sweeps, subject to the board's supply connections. \\
Keithley DMM6500 & DC voltage and current checks. \\
Keysight 34460A & DC voltage and current checks. \\
Keithley 2281S-20-6 & Programmable supply and supply-sensitivity checks. \\
Keithley 2230-30-1 & Board and bias supply rails. \\
Keysight U8031A & Board and bias supply rails. \\
Keysight E3630A & Board and bias supply rails within its available voltage and current ranges. \\
\addlinespace

\multicolumn{2}{@{}l}{\textbf{Board assembly and rework}} \\
Aoyue 932 & Vacuum pickup for component placement. \\
Weller WXR3 & Soldering and desoldering. \\
Quick 957DW+ & Hot-air rework. \\
\end{longtable}
\normalsize

The station checkout also provides general hand tools such as tweezers, wire cutters, and screwdrivers [1].

\section{Measurement Configuration and Remaining Requirements}

\begin{itemize}
    \item \textbf{RF and LO generation:}
    Conversion-gain measurements require independently set RF and LO signals near 900 MHz.
    Two-tone IIP3 measurements additionally require two simultaneous RF tones, generated by two RF sources and a suitable combiner or by verified multitone capability, while retaining an independent LO.
    The public inventory does not confirm that complete source configuration.

    \item \textbf{900 MHz network measurements:}
    E5063A option 205 stops at 500 MHz and cannot cover the nominal RF frequency.
    Option 215 or a higher-frequency option is needed; the installed option and a suitable calibration kit must be checked [2].
    A standard same-frequency S-parameter measurement is not by itself a measurement of frequency-converting gain.

    \item \textbf{IF generation and digitization:}
    The AFG31152 supports sine-wave generation to 150 MHz; the 33500B family offers 20 or 30 MHz bandwidth depending on model [3,4].
    Both are candidates for the 0.5--5 MHz IF test chain.
    Oscilloscope sample rate, record length, probe impedance, and analog bandwidth must be set for the measurement; the proposal does not specify a 2 MS/s ADC.

    \item \textbf{Analyzer and load conditions:}
    The N9000B family covers 900 MHz and the proposed IF band, but coupling, installed applications, noise floor, and calibration must be checked [5].
    Its 50-ohm input must not be treated as the specified high-impedance IF load without a characterized interface or loading correction.

    \item \textbf{Noise figure and RF accessories:}
    A calibrated noise source or equivalent calibrated noise-measurement method, suitable IF measurement chain, RF power sensor, cables, attenuators, terminations, and calibration standards must be identified.
    These accessories and a licensed NF application are not established by the benchtop list.
    Mixer NF measurements must account for image-band noise and the selected SSB definition.

    \item \textbf{Harness and environmental tests:}
    Board-specific control, supply, and pad connections must follow the final harness documentation.
    Temperature-control equipment for measured temperature sweeps is not identified in this inventory; post-layout PVT simulation remains part of the design verification plan.
\end{itemize}

\section{Sources}

\begingroup
\small
\begin{enumerate}[label={[\arabic*]},nosep]
    \item \href{https://wiki.hivemakerspace.com/explore/electronics-benchtop/}{The Hive: Electronics Benchtop inventory}, reviewed November 19, 2026.
    \item \href{https://helpfiles.keysight.com/csg/e5063a/product_information/options.htm}{Keysight: E5063A frequency options}.
    \item \href{https://www.tek.com/en/datasheet/arbitrary-function-generators}{Tektronix: AFG31000 series specifications}.
    \item \href{https://www.keysight.com/content/dam/keysight/en/doc/ungate/product-fact-sheets/5991-0757.pdf}{Keysight: 33500B series product fact sheet}.
    \item \href{https://www.keysight.com/us/en/assets/9018-04781/technical-specifications/9018-04781.pdf}{Keysight: N9000B CXA specification guide}.
    \item \href{https://opencircuitdesign.com/chipalooza/rules-4.html}{Chipalooza Challenge \#4 Rules and Regulations}, especially Rule 6.6.
\end{enumerate}
\endgroup

\end{document}
